\SourceLeaf{86}%
\Article{13}
Si exempla art. praec. inter se comparantur, primo saltem aspectu criterium nullum simplex se offerre videtur, per quod modulos priores a posterioribus dignoscere liceret. Nihilominus \emph{duo} huiusmodi criteria dantur, elegantia et simplicitate perinsignia, ad quorum alterum considerationes sequentes viam sternent.

Modulus $p$, tamquam numerus primus formae $8n+1$, reduci poterit, et quidem unico tantum modo, sub formam $aa+2bb$ (\emph{Disquiss. Arithm.} art. 182, II); radices $a,b$ positive accipi supponemus. Manifesto $a$ impar erit, $b$ vero par; statuemus autem $b=2^\lambda c$, ita ut $c$ sit impar. Iam observamus

I. quum habeatur $p\equiv aa\ (\mathrm{mod.}\ c)$, ipsum $p$ esse residuum quadraticum ipsius $c$, et proin etiam singulorum factorum primorum, in quos $c$ resolvitur: vicissim itaque, per theorema fundamentale, singuli hi factores primi erunt residua quadratica ipsius $p$, et proin etiam illorum productum $c$ erit residuum quadraticum ipsius $p$. Quod quum etiam de numero $2$ valeat, patet, $b$ esse residuum quadraticum ipsius $p$, et proin $bb$, nec non $-bb$, residuum biquadraticum.

II. Hinc $-2bb$ ad eandem classem referri debet, in qua invenitur numerus $2$; quare quum $aa\equiv-2bb$, manifestum est, $2$ vel in classe $A$, vel in classe $C$ inveniri, prout $a$ sit vel residuum quadraticum ipsius $p$, vel non-residuum quadraticum.

III. Iam supponamus, $a$ in factores suos primos resolutum esse, e quibus ii, qui sunt vel formae $8m+1$ vel $8m+7$, denotentur per $\alpha,\alpha',\alpha''$ etc., ii vero, qui sunt vel formae $8m+3$ vel $8m+5$, per $\beta,\beta',\beta''$ etc.: posteriorum multitudo sit $=\mu$. Quoniam $p\equiv2bb\ (\mathrm{mod.}\ a)$, erit $p$ residuum quadraticum eorum factorum primorum ipsius $a$, quorum residuum quadraticum est $2$, i. e. factorum $\alpha,\alpha',\alpha''$ etc.; non-residuum quadraticum vero factorum eorum, quorum non-residuum quadraticum est $2$, i. e. factorum $\beta,\beta',\beta''$ etc. Quocirca, vice versa, per theorema fundamentale, singuli $\alpha,\alpha',\alpha''$ etc. erunt residua quadratica ipsius $p$, singuli $\beta,\beta',\beta''$ etc. autem non-residua quadratica. Ex his itaque concluditur, productum $a$ fore residuum quadraticum ipsius $p$, vel non-residuum, prout $\mu$ par sit vel impar.

IV. Sed facile confirmatur, productum omnium $\alpha,\alpha',\alpha''$ etc. fieri formae $8m+1$ vel $8m+7$, idemque valere de producto omnium $\beta,\beta',\beta''$ etc., si horum multitudo fuerit par, ita ut in hoc casu etiam productum $a$ necessario fieri debeat formae $8m+1$ vel $8m+7$; contra productum omnium $\beta,\beta',\beta''$ etc., quo%
\SourceLeaf{87}%
ties ipsorum multitudo impar sit, fieri formae $8m+3$ vel $8m+5$, idemque adeo in hoc casu valere de producto $a$.

Ex his omnibus itaque colligitur theorema elegans:

\emph{Quoties $a$ est formae $8m+1$ vel $8m+7$, numerus $2$ in complexu $A$ contentus erit; quoties vero $a$ est formae $8m+3$ vel $8m+5$, numerus $2$ in complexu $C$ invenietur.}

Quod confirmatur per exempla in art. praec. enumerata; priores enim moduli ita discerpuntur: $73=1+2\pdot36$, $89=81+2\pdot4$, $113=81+2\pdot16$, $233=225+2\pdot4$, $257=225+2\pdot16$, $281=81+2\pdot100$, $337=49+2\pdot144$, $353=225+2\pdot64$; posteriores vero ita: $17=9+2\pdot4$, $41=9+2\pdot16$, $97=25+2\pdot36$, $137=9+2\pdot64$, $193=121+2\pdot36$, $241=169+2\pdot36$, $313=25+2\pdot144$, $401=9+2\pdot196$, $409=121+2\pdot144$, $433=361+2\pdot36$, $449=441+2\pdot4$, $457=169+2\pdot144$.

\Article{14}
Quum discerptio numeri $p$ in quadratum simplex et duplex nexum tam insignem cum classificatione numeri $2$ prodiderit, operae pretium esse videtur tentare, num discerptio in duo quadrata, cui numerum $p$ aeque obnoxium esse constat, similem forte successum suppeditet. Ecce itaque discerptiones numerorum $p$, pro quibus $2$ pertinet ad classem
\[
\begin{array}{c|c}
A&C\\\hline
9+64&1+16\\
25+64&25+16\\
49+64&81+16\\
169+64&121+16\\
1+256&49+144\\
25+256&225+16\\
81+256&169+144\\
289+64&1+400\\
&9+400\\
&289+144\\
&49+400\\
&441+16
\end{array}
\]

\SourceLeaf{88}%
Ante omnia observamus, duorum quadratorum . in quae $p$ discerpitur, alterum impar esse debere, quod statuemus $=aa$, alterum par, quod statuemus $=bb$. Quoniam $aa$ fit formae $8n+1$, patet, valoribus impariter paribus ipsius $b$ respondere valores ipsius $p$ formae $8n+5$, ab inductione nostra hic exclusos, quippe qui numerum $2$ in classe $B$ vel $D$ haberent. Pro valoribus autem ipsius $p$, qui sunt formae $8n+1$, $b$ esse debet pariter par, et si inductioni, quam schema allatum ob oculos sistit, fidem habere licet, numerus $2$ ad classem $A$ referendus erit pro omnibus modulis, pro quibus $b$ est formae $8n$, ad classem $C$ vero pro omnibus modulis, pro quibus $b$ est formae $8n+4$. Sed hoc theorema longe altioris indaginis est, quam id, quod in art. praec. eruimus, demonstrationique plures disquisitiones praeliminares sunt praemittendae, ordinem, quo numeri complexum $A,B,C,D$ se invicem sequuntur, spectantes.

\Article{15}
Designemus multitudinem numerorum e complexu $A$, quos immediate sequitur numerus e complexu $A,B,C,D$ resp., per $(00),(01),(02),(03)$; perinde multitudinem numerorum e complexu $B$, quos sequitur numerus e complexu $A,B,C,D$ resp. per $(10),(11),(12),(13)$; similiterque sint in complexu $C$ resp. $(20),(21),(22),(23)$ numeri, in complexu $D$ vero $(30),(31),(32),(33)$ numeri quos sequitur numerus e complexu $A,B,C,D$. Proponimus nobis, has sedecim multitudines a priori determinare. Quo commodius lectores ratiocinia generalia cum exemplis comparare possint. valores numericos terminorum schematis $(S)$
\[
\begin{array}{llll}
(00),&(01),&(02),&(03)\\
(10),&(11),&(12),&(13)\\
(20).&(21),&(22),&(23)\\
(30),&(31),&(32),&(33)
\end{array}
\]
pro singulis modulis, pro quibus classificationes in art. 11 tradidimus, hic adscribere visum est.

\[
\begin{array}{c|c|c|c}
\begin{gathered}
p=5\\[3pt]
\begin{array}{l}
0,1,0\mathpunct{.}\,0\\
0,0,0,1\\
0,0,0,0\\
0,0,1,0
\end{array}
\end{gathered}&
\begin{gathered}
p=13\\[3pt]
\begin{array}{l}
0,1,2,0\\
1,1,0,1\\
0,1,0,1\\
1,0,1,1
\end{array}
\end{gathered}&
\begin{gathered}
p=17\\[3pt]
\begin{array}{l}
0,2,1,0\\
2,0,1,1\\
1,1,1,1\\
0,1,1,2
\end{array}
\end{gathered}&
\begin{gathered}
p=29\\[3pt]
\begin{array}{l}
2\mathpunct{.}\,3,0\mathpunct{.}\,2\\
1,1,2,3\\
2,1,2\mathpunct{.}\,1\\
1,2,3,1
\end{array}
\end{gathered}
\end{array}
\]

\SourceLeaf{89}%
\[
\begin{array}{c|c|c|c}
\begin{gathered}
p=37\\[3pt]
\begin{array}{l}
2,1,2,4\\
2,2,4,1\\
2,2,2,2\\
2,4,1,2
\end{array}
\end{gathered}&
\begin{gathered}
p=41\\[3pt]
\begin{array}{l}
0,4,3,2\\
4,2,2,2\\
3,2,3,2\\
2,2,2,4
\end{array}
\end{gathered}&
\begin{gathered}
p=53\\[3pt]
\begin{array}{l}
2,3,6,2\\
4,4,2,3\\
2,4,2,4\\
4,2,3,4
\end{array}
\end{gathered}&
\begin{gathered}
p=61\\[3pt]
\begin{array}{l}
4,3,2,6\\
3,3,6,3\\
4,3,4,3\\
3,6,3,3
\end{array}
\end{gathered}
\end{array}
\]

\[
\begin{array}{c|c|c}
\begin{gathered}
p=73\\[3pt]
\begin{array}{l}
5,6,4,2\\
6,2,5,5\\
4,5,4,5\\
2,5,5,6
\end{array}
\end{gathered}&
\begin{gathered}
p=89\\[3pt]
\begin{array}{l}
3,8,6,4\\
8,4,5,5\\
6,5,6,5\\
4,5,5,8
\end{array}
\end{gathered}&
\begin{gathered}
p=97\\[3pt]
\begin{array}{l}
2,6,7,8\\
6,8,5,5\\
7,5,7,5\\
8,5,5,6
\end{array}
\end{gathered}
\end{array}
\]

Quum moduli formae $8n+1$ et $8n+5$ diverso modo se habeant, utrosque seorsim tractare oportet: a prioribus initium faciemus.

\Article{16}
Character $(00)$ indicat, quot modis diversis aequationi $\alpha+1=\alpha'$ satisfieri possit, denotantibus $\alpha,\alpha'$ indefinite numeros e complexu $A$. Quum pro modulo formae $8n+1$, qualem hic subintelligimus, $\alpha'$ et $p-\alpha'$ ad eundem complexum pertineant, concinnius dicemus, $(00)$ exprimere multitudinem modorum diversorum, aequationi $1+\alpha+\alpha'=p$, satisfaciendi: manifesto huius aequationis vice etiam congruentia $1+\alpha+\alpha'\equiv0\ (\mathrm{mod.}\ p$ fungi potest.

Perinde
\begin{center}
\begin{minipage}{\linewidth}
$(01)$ indicat multitudinem solutionum congruentiae $1+\alpha+\beta\equiv0\ (\mathrm{mod.}\ p)$\par
$(02)$ multitudinem solutionum congruentiae $1+\alpha+\gamma\equiv0$\par
$(03)$ multitudinem solutionum congruentiae $1+\alpha+\delta\equiv0$\par
$(11)$ multitudinem solutionum congruentiae $1+\beta+\beta'\equiv0$ etc.
\end{minipage}
\end{center}
exprimendo indefinite per $\beta$ et $\beta'$ numeros e complexu $B$, per $\gamma$ numeros e complexu $C$, per $\delta$ numeros e complexu $D$. Hinc statim colligimus sex aequationes sequentes:
\[
\begin{gathered}
(01)=(10),\quad(02)=(20),\quad(03)=(30),\\
(12)=(21),\quad(13)=(31),\quad(23)=(32)
\end{gathered}
\]

E quavis solutione data congruentiae $1+\alpha+\beta\equiv0$ demanat solutio congruentiae $1+\delta+\delta'\equiv0$, accipiendo pro $\delta$ numerum inter limites $1\dotsfour p-1$
\SourceLeaf{90}%
eum qui reddit $\beta\delta\equiv1$ (qui manifesto erit e complexu $D$), et pro $\delta'$ residuum minimum positivum producti $\alpha\delta$ (quod itidem erit e complexu $D$); perinde patet regressus a solutione data congruentiae $1+\delta+\delta'\equiv0$ ad solutionem congruentiae $1+\alpha+\beta\equiv0$, si $\beta$ accipitur ita, ut fiat $\beta\delta\equiv1$, simulque statuitur $\alpha\equiv\beta\delta'$. Hinc concludimus, utramque congruentiam aequali solutionum multitudine gaudere, sive esse $(01)=(33)$.

Simili modo e congruentia $1+\alpha+\gamma\equiv0$ deducimus $\gamma'+\gamma''+1\equiv0$, si $\gamma'$ accipitur e complexu $C$ ita ut fiat $\gamma\gamma'\equiv1$, atque $\gamma''$ ex eodem complexu congruus producto $\alpha\gamma'$. Unde facile colligimus, has duas congruentias aequalem solutionum multitudinem admittere, sive esse $(02)=(22)$.

Perinde e congruentia $1+\alpha+\delta\equiv0$ deducimus $\beta+\beta'+1\equiv0$, accipiendo $\beta,\beta'$ ita ut fiat $\beta\delta\equiv1$, $\beta\alpha\equiv\beta'$, eritque adeo $(03)=(11)$.

Denique e congruentia $1+\beta+\gamma\equiv0$ simili modo tum congruentiam $\delta+1+\beta'\equiv0$, tum hanc $\gamma'+\delta'+1\equiv0$ derivamus, atque hinc concludimus $(12)=(13)=(23)$.

Nacti sumus itaque, inter sedecim incognitas nostras, undecim aequationes, ita ut illae ad quinque reducantur, schemaque $S$ ita exhiberi possit:
\[
\begin{array}{llll}
h,&i,&k,&l\\
i,&l,&m,&m\\
k,&m,&k,&m\\
l,&m,&m,&i
\end{array}
\]

Facile vero tres novae aequationes conditionales adiiciuntur. Quum enim quemvis numerum complexus $A$, excepto ultimo $p-1$, sequi debeat numerus ex aliquo complexuum $A,B,C$ vel $D$, habebimus
\[
(00)+(01)+(02)+(03)=2n-1
\]
et perinde
\[
\begin{aligned}
(10)+(11)+(12)+(13)&=2n\\
(20)+(21)+(22)+(23)&=2n\\
(30)+(31)+(32)+(33)&=2n
\end{aligned}
\]
In signis modo introductis tres primae aequationes suppeditant:
\[
\begin{aligned}
h+i+k+l&=2n-1\\
i+l+2m&=2n\\
k+m&=n
\end{aligned}
\]

